\documentclass[10pt]{article}
\usepackage[T1]{fontenc}
\usepackage[utf8]{inputenc}
\usepackage[french]{babel}
\usepackage{amsmath,amssymb,mathrsfs,mathpazo}
\usepackage{geometry}
\geometry{a4paper,margin=2cm}
\usepackage{graphicx,array,tabularx,tabularray,colortbl,multirow,multicol,booktabs,enumitem,pifont,ulem,fancybox,diagbox,cancel}
\usepackage{pstricks,pstricks-add,pst-plot,pst-tree,pst-node,pst-text,pst-func,pst-math,pst-bspline,pst-3dplot,pst-eucl,pst-solides3d}
\let\pscalculate\undefined
\usepackage{pst-bezier}
\usepackage{tikz,pgfplots,tkz-tab}
\usetikzlibrary{arrows,arrows.meta,patterns,calc,positioning,trees,decorations.markings,shapes.geometric,shapes.misc,shapes.arrows,intersections,decorations.pathmorphing,decorations.pathreplacing,backgrounds,shadows,babel}
\pgfplotsset{compat=1.18}
\usepgfplotslibrary{fillbetween}
\usepackage[np]{numprint}
\usepackage{listings,makecell,pas-tableur,twoopt}
\let\barre\undefined
\usepackage[pstricks]{tablvar}
\definecolor{darkgreen}{RGB}{0,100,0}
\providecommand{\red}{\color{red}}
\providecommand{\blue}{\color{blue}}
\providecommand{\green}{\color{green}}
\providecommand{\black}{\color{black}}
\providecommand{\white}{\color{white}}
\providecommand{\gray}{\color{gray}}
\providecommand{\cyan}{\color{cyan}}
\providecommand{\magenta}{\color{magenta}}
\providecommand{\yellow}{\color{yellow}}
\providecommand{\violet}{\color{violet}}
\providecommand{\orange}{\color{orange}}
\providecommand{\euro}{\texteuro}
\providecommand{\EUR}{\texteuro}
\providecommand{\ieme}{\textsuperscript{e}}
\providecommand{\degres}{\ensuremath{^\circ}}
\providecommand{\degree}{\ensuremath{^\circ}}
\providecommand{\vv}[1]{\overrightarrow{#1}}
\providecommand{\Coord}[1]{#1}
\providecommand{\decosix}{}
\providecommand{\starredbullet}{\ensuremath{\bullet}}
\pagestyle{empty}
\setlength{\parindent}{0pt}
\usepackage{adjustbox}\DeclareUnicodeCharacter{2236}{:}\DeclareUnicodeCharacter{2190}{\ensuremath{\leftarrow}}
\providecommand{\R}{\mathbb{R}}
\providecommand{\N}{\mathbb{N}}
\providecommand{\Z}{\mathbb{Z}}
\providecommand{\Q}{\mathbb{Q}}
\providecommand{\D}{\mathbb{D}}
\providecommand{\C}{\mathbb{C}}
\providecommand{\e}{{\mathrm{e}}}
\providecommand{\vect}[1]{\overrightarrow{#1}}
\providecommand{\vectt}[1]{\overrightarrow{\mathrm{#1}}}
\providecommand{\vvt}[1]{\overrightarrow{\mathrm{#1}}}
\providecommand{\np}[2][]{{#2~#1}}
\providecommand{\vv}[1]{{\overrightarrow{#1}}}
\renewcommand{\d}{{\mathrm{d}}}
\providecommand{\degres}{{^\circ}}
\providecommand{\degre}{{^\circ}}
\providecommand{\blue}{\color{blue}}
\providecommand{\red}{\color{red}}
\providecommand{\enskip}{\quad}
\providecommand{\llbracket}{\lbrack\!\lbrack}
\providecommand{\rrbracket}{\rbrack\!\rbrack}
\providecommand{\euro}{{\text{€}}}
\providecommand{\no}{{\text{n°}}}
\providecommand{\interval}[2]{{[#1;#2]}}
\providecommand{\intervFO}[2]{{[#1;#2[}}
\providecommand{\dotfill}{\ldots}

\providecommand{\ds}{\displaystyle}
\providecommand{\textding}[1]{#1}
\providecommand{\Oij}{\(\mathrm{O};\vec{\imath},\vec{\jmath}\)}
\providecommand{\Oijk}{\(\mathrm{O};\vec{\imath},\vec{\jmath},\vec{k}\)}
\providecommand{\Ouv}{\(\mathrm{O};\vec{u},\vec{v}\)}
\providecommand{\intervFF}[2]{\ensuremath{[#1;#2]}}
\renewcommand{\euro}{\ifmmode\text{\texteuro}\else\texteuro\fi}
\DeclareUnicodeCharacter{00B0}{\ifmmode^{\circ}\else\textdegree\fi}
\DeclareUnicodeCharacter{20AC}{\ifmmode\text{\texteuro}\else\texteuro\fi}
% Énoncés : Éducation nationale. Transcription et corrigés : APMEP, Denis Vergès et contributeurs. Découpage et classement : CoopMaths, Sébastien Lozano, Éric Elter et contributeurs. Adaptation de la présentation : Maths Explorer.
% CC BY-SA : https://coopmaths.fr/www/about/
% https://coopmaths.fr/alea/static/e3c/2021/tex/e3c_2021_01_specimen1_2_cor.tex
\begin{document}
\noindent\textit{APMEP / CoopMaths. CC BY-SA. Présentation adaptée par Maths Explorer.}\par\medskip
\bigskip

Une urne contient six jetons rouges dont un est marqué « gagnant » et quatre jetons verts dont trois d'entre eux sont marqués « gagnant ».

\begin{list}{}{On tire au hasard un jeton de l'urne et on note les évènements :}
\item $R$ : « le jeton tiré est rouge »,
\item $V$ : « le jeton tiré est vert »,
\item $G$ : « le jeton tiré est gagnant ».
\end{list}

\begin{enumerate}
\item  On modélise la situation à l'aide d'un arbre de probabilité.

\begin{center}
\bigskip
  \pstree[treemode=R,nodesepA=0pt,nodesepB=4pt,levelsep=2.5cm,treesep=1.2cm]{\TR{}}
 {
 	\pstree[nodesepA=4pt]{\TR{$R$}\ncput*{$\frac{6}{10}$}}
 	  { 
 		  \TR{$G$}\ncput*{$\frac{1}{6}$}
 		  \TR{$\overline{G}$}\ncput*{$\frac{5}{6}$}	   
 	  }
 	\pstree[nodesepA=4pt]{\TR{$V$}\ncput*{$\frac{4}{10}$}}
 	  {
 		  \TR{$G$}\ncput*{$\frac{3}{4}$}
          \TR{$\overline{G}$}\ncput*{$\frac{1}{4}$} 
     }
}
\bigskip
\end{center}

\item La probabilité de l'évènement « le jeton tiré est un jeton vert et marqué gagnant » est:

$P(V\cap G) = p(V)\times P_{V}(G)= \dfrac{4}{10}\times \dfrac{3}{4}=\dfrac{3}{10}$.

\item Soit $P(G)$ la probabilité de tirer un jeton gagnant. 

D'après la formule des probabilités totales:

$P(G) = P(R\cap G) + P(V\cap G) = \dfrac{6}{10}\times \dfrac{1}{6}+ \dfrac{3}{10} 
= \dfrac{1}{10} + \dfrac{3}{10} = \dfrac{4}{10}=\dfrac{2}{5}$.

\item Sachant que le jeton tiré est gagnant, la probabilité qu'il soit de couleur rouge est:

$P_G(R) = \dfrac{P(R\cap G)}{P(G)} = \dfrac{\frac{1}{10}}{\frac{4}{10}} = \dfrac{1}{4}$.

\item On tire maintenant, toujours au hasard et simultanément, deux jetons dans l'urne.

On cherche la probabilité que les deux jetons soient marqués « gagnant ». 



Il y a en tout 10 jetons dont 4 gagnants.

\begin{list}{\textbullet}{}
\item Il y a $\ds\binom{10}{2}=45$ façons de tirer 2 jetons parmi 10.
\item Il y a $\ds\binom{4}{2}= 6$ façons de tirer 2 jetons gagnants parmi les 4 gagnants.
\item Le tirage se fait au hasard donc il y a équiprobabilité.

\end{list}

La probabilité que les deux jetons soient marqués « gagnant » est donc: $\dfrac{6}{45} = \dfrac{2}{15}$. 

\end{enumerate}

\end{document}
